\documentclass[aspectratio=169,14pt]{beamer}
\usepackage[T1]{fontenc}
\usepackage[utf8]{inputenc}
\usepackage[french]{babel}
\usepackage{amsthm}
\usepackage{amsfonts}
\usepackage{amsmath}
\usepackage{amssymb}
\usepackage{mathtools}
\usepackage{bbm}
\usepackage{pgfplots}
\usepackage{tikz}
\usepackage{physics}
\usepackage{calligra}
\usepackage{csquotes}
\usepackage{tensor}
\usepackage[thicklines]{cancel}
\usepackage{tcolorbox}
\usepackage{pstricks}
\usepackage[backend=biber, bibstyle=nature, sorting=nty, citestyle=numeric-comp]{biblatex} %Custom bibliography
    \addbibresource{bib.bib} %Load references


\DeclareMathAlphabet{\mathcalligra}{T1}{calligra}{m}{n}
\DeclareFontShape{T1}{calligra}{m}{n}{<->s*[2.2]callig15}{}
\newcommand{\scriptr}{\mathcalligra{r}\,}
\newcommand{\boldscriptr}{\pmb{\mathcalligra{r}}\,}
\def\rc{\scriptr}
\def\brc{\boldscriptr}
\def\hrc{\hat\brc}
\newcommand{\ie}{\emph{i.e.}} %id est
\newcommand{\eg}{\emph{e.g.}} %exempli gratia
\newcommand{\rtd}[1]{\ensuremath{\left\lfloor #1 \right\rfloor}}
\newcommand{\dirac}[1]{\ensuremath{\delta \left( #1 \right)}}
\newcommand{\diract}[1]{\ensuremath{\delta^3 \left( #1 \right)}}
\newcommand{\e}{\ensuremath{\epsilon_0}}
\newcommand{\m}{\ensuremath{\mu_0}}
\newcommand{\V}{\ensuremath{\mathcal{V}}}
\newcommand{\prnt}[1]{\ensuremath{\left(#1\right)}} %parentheses
\newcommand{\colch}[1]{\ensuremath{\left[#1\right]}} %square brackets
\newcommand{\chave}[1]{\ensuremath{\left\{#1\right\}}}  %curly brackets
\newcommand\eqdef{\stackrel{\mathclap{\normalfont \tiny\mbox{\textrm{def}}}}{=}}


\definecolor{blue2}{HTML}{045FB4}
\definecolor{green2}{HTML}{46C235}
\definecolor{red2}{HTML}{EE4848}
\definecolor{violet2}{HTML}{A647E5}
\definecolor{orange2}{HTML}{FF7425}
\definecolor{darkred}{HTML}{5C2020}
\definecolor{gray}{HTML}{303030}
\definecolor{yellow}{HTML}{f0be52}
\definecolor{lightdarkgold}{HTML}{EEBC1D}

\renewcommand{\CancelColor}{\color{darkred}}

\makeatletter
\newcommand{\mybox}[1]{%
  \setbox0=\hbox{#1}%
  \setlength{\@tempdima}{\dimexpr\wd0+13pt}%
  \begin{tcolorbox}[colback=darkred,colframe=darkred,boxrule=0.5pt,arc=4pt,
      left=6pt,right=6pt,top=6pt,bottom=6pt,boxsep=0pt,width=\@tempdima]
    \textcolor{yellow}{#1}
  \end{tcolorbox}
}
\makeatother


\pgfplotsset{my style/.append style={axis x line=middle, axis y line=
middle, xlabel={$x$}, ylabel={$y$}, axis equal }}


% Le thème visuel est dans chapters/theme.tex

\newcommand{\hlight}[1]{\colorbox{violet!50}{#1}}
\newcommand{\hlighta}[1]{\colorbox{darkred!50}{#1}}
\usetikzlibrary{positioning,calc,arrows.meta}

% Petite macro : source / référence en bas de slide
\newcommand{\src}[1]{\par\vfill{\scriptsize\textcolor{gray!70}{#1}}}

\title{IA et neurosciences du langage}
\subtitle{Ce que les modèles nous apprennent sur le cerveau, et inversement}
\author[B. Sow]{\textcolor{darkred}{Boubacar Sow}}
\institute[LSCP]{
    \textcolor{gray}{Doctorant, Laboratoire de Sciences Cognitives et Psycholinguistique}%
    \\%
    \textcolor{gray}{ENS-PSL \& CNRS}%
}
\titlegraphic{%
    \includegraphics[height=1.1cm]{images/lscp.png}\hspace{1.2cm}%
    \includegraphics[height=1.4cm]{images/ens_psl.png}\hspace{1.2cm}%
    \includegraphics[height=1.1cm]{images/cnrs.png}%
}
\date{IndabaX Guinée 2026} % TODO : date exacte
%\logo{\includegraphics[width= 0.05\textwidth]{images/logo.png}}

\begin{document}

    \frame{\titlepage}

    \begin{frame}{Plan}
        \tableofcontents
    \end{frame}

    \section{Problématique et cadre théorique}

\begin{frame}{Problématique}
\begin{columns}[T]
    \column{.48\textwidth}
    \textcolor{navy}{\textbf{Cerveau humain}}
    \begin{itemize}
        \item $\sim$86 milliards de neurones
        \item $\sim$20 W
        \item Acquisition du langage en quelques années d'exposition
    \end{itemize}
    \column{.48\textwidth}
    \textcolor{navy}{\textbf{Modèle de langage}}
    \begin{itemize}
        \item $10^9$ à $10^{11}$ paramètres
        \item Entraînement sur $\sim 10^{12}$ mots
        \item Performance proche de l'humain en production
    \end{itemize}
\end{columns}
\vspace{1em}
\begin{block}{Question}
    Ces deux systèmes mettent-ils en œuvre des calculs comparables ?
\end{block}
\end{frame}

\begin{frame}{Précédents historiques}
\begin{itemize}
    \item 1943 : neurone formel (McCulloch \& Pitts) ; 1958 : perceptron (Rosenblatt)
    \item Réseaux convolutifs inspirés du cortex visuel (Hubel \& Wiesel)
    \item 2014 : des réseaux optimisés pour la reconnaissance d'objets prédisent l'activité du cortex visuel des primates
\end{itemize}
\vspace{0.5em}
\begin{block}{Extension au langage}
    Question abordée depuis environ 2018 avec les modèles de langage.
\end{block}
\src{Yamins et al., PNAS 2014}
\end{frame}

\begin{frame}{Les trois niveaux de Marr}
\begin{columns}[T]
    \column{.52\textwidth}
    \begin{itemize}
        \item<1-> \textbf{Computationnel} : \emph{quel problème} est résolu ?\\
              {\small ex. : une addition}
        \item<2-> \textbf{Algorithmique} : \emph{quelles représentations et quelles procédures} ?\\
              {\small ex. : chiffres arabes et retenue}
        \item<3-> \textbf{Implémentationnel} : \emph{substrat} physique\\
              {\small ex. : neurones, transistors}
    \end{itemize}
    \column{.46\textwidth}
    \centering
    \resizebox{\linewidth}{!}{%
    \begin{tikzpicture}[
        n/.style={circle, draw=navy, fill=navy!15, minimum size=0.42cm, inner sep=0pt},
        e/.style={navy, thick},
        lab/.style={font=\small\bfseries, text=navy, anchor=west}]
        \node<1->[n] (c) at (0,0) {};
        \node<1->[lab] at (2.4,0) {Computationnel};
        \node<2->[n] (a1) at (-1.2,-1.4) {};
        \node<2->[n] (a2) at (0,-1.4) {};
        \node<2->[n] (a3) at (1.2,-1.4) {};
        \node<2->[lab] at (2.4,-1.4) {Algorithmique};
        \draw<2->[e] (c)--(a1); \draw<2->[e] (c)--(a2); \draw<2->[e] (c)--(a3);
        \foreach \x/\i in {-1.8/1,-0.6/2,0.6/3,1.8/4}{
            \node<3->[n] (i\i) at (\x,-2.8) {};
            \draw<3->[e] (a2)--(i\i);}
        \node<3->[lab] at (2.4,-2.8) {Implémentation};
        \node<3->[font=\small\bfseries, text=navy] at (0,-3.5) {Relation un-à-plusieurs};
    \end{tikzpicture}}
\end{columns}
\end{frame}

\begin{frame}{Marr : humains et modèles de langage}
\centering
\resizebox{\textwidth}{!}{%
\begin{tikzpicture}[
    n/.style={circle, draw=navy, fill=navy!15, minimum size=0.42cm, inner sep=0pt},
    e/.style={navy, thick},
    lab/.style={font=\small\bfseries, text=navy, anchor=west}]
    \node[n] (c) at (0,0) {};
    \foreach \x/\i in {-4/1,-2/2,0/3,2/4,4/5}{\node[n] (a\i) at (\x,-1.5) {};}
    \foreach \i in {2,3,4}{\draw[e] (c)--(a\i);}
    \draw[lscp, ultra thick] (c)--(a1);
    \draw[blue2, ultra thick] (c)--(a5);
    \foreach \x/\i in {-5.2/1,-4.4/2,-3.6/3,-2.8/4}{\node[n] (m\i) at (\x,-3.2) {}; \draw[e] (a1)--(m\i);}
    \foreach \x/\i in {2.8/1,3.6/2,4.4/3,5.2/4}{\node[n] (h\i) at (\x,-3.2) {}; \draw[e] (a5)--(h\i);}
    \node[lab, text=lscp!70!black] at (-5.2,-3.9) {Modèles de langage};
    \node[lab, text=blue2] at (2.6,-3.9) {Cerveaux humains};
    \node[lab] at (6.2,0) {Computationnel};
    \node[lab] at (6.2,-1.5) {Algorithmique};
    \node[lab] at (6.2,-3.2) {Implémentation};
\end{tikzpicture}}

\vspace{0.4em}
\textbf{Hypothèse de divergence} : implémentations très différentes, éventuellement objectifs différents ; les algorithmes diffèrent.
\end{frame}

\begin{frame}{Marr : humains et modèles de langage}
\centering
\resizebox{\textwidth}{!}{%
\begin{tikzpicture}[
    n/.style={circle, draw=navy, fill=navy!15, minimum size=0.42cm, inner sep=0pt},
    e/.style={navy, thick},
    lab/.style={font=\small\bfseries, text=navy, anchor=west}]
    \node[n] (c) at (0,0) {};
    \foreach \x/\i in {-4/1,-2/2,0/3,2/4,4/5}{\node[n] (a\i) at (\x,-1.5) {};}
    \foreach \i in {1,2,4,5}{\draw[e] (c)--(a\i);}
    \draw[lscp, ultra thick] (c) to[bend right=28] (a3);
    \draw[blue2, ultra thick] (c) to[bend left=28] (a3);
    \foreach \x/\i in {-5.2/1,-4.4/2,-3.6/3,-2.8/4}{\node[n] (m\i) at (\x,-3.2) {}; \draw[e] (a3)--(m\i);}
    \foreach \x/\i in {2.8/1,3.6/2,4.4/3,5.2/4}{\node[n] (h\i) at (\x,-3.2) {}; \draw[e] (a3)--(h\i);}
    \node[lab, text=lscp!70!black] at (-5.2,-3.9) {Modèles de langage};
    \node[lab, text=blue2] at (2.6,-3.9) {Cerveaux humains};
    \node[lab] at (6.2,0) {Computationnel};
    \node[lab] at (6.2,-1.5) {Algorithmique};
    \node[lab] at (6.2,-3.2) {Implémentation};
\end{tikzpicture}}

\vspace{0.4em}
\textbf{Hypothèse de convergence} : un \emph{même} algorithme, implémenté sur deux substrats différents.
\end{frame}

\begin{frame}{Inférer le niveau algorithmique}
\begin{itemize}
    \item L'algorithme n'est pas directement observable : il est \textbf{inféré}
    \item \textbf{Comportement} : temps de lecture (lecture auto-rythmée, oculométrie)
    \item \textbf{Activité cérébrale} : motifs d'activité reflétant les représentations utilisées
    \item On compare donc modèles et humains \textbf{au niveau algorithmique}
\end{itemize}
\vspace{0.4em}
\begin{block}{Intérêt des modèles d'IA}
    Sous l'hypothèse de convergence, un modèle est une hypothèse explicite sur l'algorithme cérébral, entièrement observable et manipulable (ablation, intervention), ce que le cerveau permet peu.
\end{block}
\src{Marr, Vision, 1982}
\end{frame}

\begin{frame}{Domaines abordés}
\begin{columns}[T]
    \column{.48\textwidth}
    \begin{block}{Modèle du système cognitif}
        \small
        Tester des hypothèses sur le calcul cérébral\\[0.3em]
        Vision et récompense ; langage (alignement) ; comportement ; mécanismes
    \end{block}
    \column{.48\textwidth}
    \begin{block}{Source d'inspiration}
        \small
        Le cerveau guide la conception des modèles\\[0.3em]
        Signal cérébral comme supervision ; apprentissage à échelle développementale
    \end{block}
\end{columns}
\vspace{0.8em}
Pour chaque domaine : \textbf{potentiel} et \textbf{état actuel} des résultats.
\end{frame}

    \section{Les bases, sans prérequis}
\frame{\sectionpage}

\begin{frame}{Le langage dans le cerveau}
\begin{itemize}
    \item Un \textbf{réseau du langage} fronto-temporal, principalement dans l'hémisphère gauche
    \item Il répond à la langue, écrite ou parlée, et peu à autre chose
    \item Il est distinct du réseau de la pensée abstraite et du raisonnement
\end{itemize}
% IMAGE : cerveau avec le réseau du langage (à ajouter)
\src{Fedorenko et al. ; Fedorenko, Open Encyclopedia of Cognitive Science, 2025}
\end{frame}

\begin{frame}{Comment « écouter » le cerveau ?}
\centering\small
\begin{tabular}{lcccc}
    \textcolor{yellow}{Méthode} & \textcolor{yellow}{Espace} & \textcolor{yellow}{Temps} & \textcolor{yellow}{Invasif ?} \\
    \hline
    IRM fonctionnelle (fMRI) & millimètres & secondes & non \\
    EEG / MEG & centimètres & millisecondes & non \\
    Intracrânien (iEEG, sEEG) & millimètres--cm & millisecondes & oui \\
    Neurones uniques & micromètres & millisecondes & oui \\
\end{tabular}
\vspace{1em}

\normalsize
L'intracrânien se fait chez des patients épileptiques implantés pour raisons cliniques.
% IMAGE : schémas des méthodes (à ajouter)
\end{frame}

\begin{frame}{Le protocole type}
\begin{enumerate}
    \item Un participant lit ou écoute un texte (souvent une histoire naturelle)
    \item On enregistre son activité cérébrale mot par mot
    \item On donne le même texte à un modèle de langage
    \item On compare les deux
\end{enumerate}
\end{frame}

\begin{frame}{Comparer un modèle à un cerveau}
\centering
\begin{tikzpicture}[
    box/.style={draw=yellow, rounded corners, fill=darkred, text=white, align=center, minimum height=1.1cm, font=\small},
    arr/.style={-{Stealth}, thick, yellow}]
    \node[box] (txt) at (0,0) {Texte\\{\footnotesize « Le chat dort… »}};
    \node[box] (lm) at (4.6,1.1) {Modèle de langage\\{\footnotesize activations}};
    \node[box] (br) at (4.6,-1.1) {Cerveau\\{\footnotesize fMRI / EEG / iEEG}};
    \node[box] (reg) at (9.2,0) {Régression\\linéaire};
    \node[box] (sc) at (12.2,0) {Score\\d'alignement};
    \draw[arr] (txt) -- (lm);
    \draw[arr] (txt) -- (br);
    \draw[arr] (lm) -- (reg);
    \draw[arr] (br) -- (reg);
    \draw[arr] (reg) -- (sc);
\end{tikzpicture}

\vspace{1em}
\begin{itemize}
    \item On prédit le signal cérébral à partir des activations du modèle
    \item Le score est mesuré sur des données \textbf{non vues} pendant l'ajustement
    \item On le compare à un \textbf{plafond de bruit} (la variabilité entre participants)
\end{itemize}
\end{frame}

\begin{frame}{Deux outils à retenir}
\begin{block}{Modèle d'encodage}
    Régression des activations du modèle vers le signal cérébral (utilisé dans Brain-Score).
\end{block}
\begin{block}{Analyse de similarité représentationnelle (RSA)}
    On compare les distances entre stimuli dans le modèle à celles mesurées dans le cerveau.
\end{block}
\end{frame}

    \section{Q1 : modèles et cerveau calculent-ils pareil ?}

\begin{frame}{Les modèles prédisent le réseau du langage}
\begin{itemize}
    \item Des dizaines de modèles de langage comparés à des données fMRI, ECoG et EEG/MEG
    \item Les meilleurs modèles sont ceux qui prédisent le mieux le mot suivant
    \item Certains approchent le plafond de bruit sur des phrases
\end{itemize}
\vspace{0.5em}
\begin{block}{Idée clé}
    Prédire le prochain mot suffit à produire des représentations proches de celles du cerveau.
\end{block}
% TODO : vérifier les chiffres exacts avant d'ajouter (Schrimpf 2021)
\src{Schrimpf et al., PNAS 2021 ; Caucheteux \& King, 2022}
\end{frame}

\begin{frame}{Le cerveau prédit aussi}
Enregistrements intracrâniens (ECoG) pendant une conversation naturelle :
\begin{itemize}
    \item Le cerveau \textbf{anticipe} le mot suivant avant qu'il soit prononcé
    \item Il réagit à la \textbf{surprise} après le mot
    \item Les mots sont représentés par des plongements contextuels, comme dans un modèle
\end{itemize}
% IMAGE : figure d'encodage temporel (à ajouter, avec autorisation de la figure)
\src{Goldstein et al., Nature Neuroscience 2022}
\end{frame}

\begin{frame}{Des données naturelles, à grande échelle}
\begin{itemize}
    \item \textbf{Narratives} : fMRI de centaines de participants qui écoutent des histoires
    \item Permet de séparer ce qui est lexical, syntaxique ou compositionnel
    \item Résultat : syntaxe et sémantique partagent un substrat distribué plutôt que des modules séparés
\end{itemize}
% TODO : vérifier la citation exacte du jeu de données (Nastase et al., 2021)
\src{Nastase et al., Narratives ; Caucheteux, Gramfort \& King, ICML 2021}
\end{frame}

\begin{frame}{Combien de données faut-il ?}
\begin{itemize}
    \item Critique fréquente : les modèles voient bien plus de texte qu'un enfant
    \item Réponse : environ \textbf{100 millions de mots} suffisent pour prédire presque au maximum les réponses fMRI à des phrases
    \item Un enfant en entend environ ce volume au cours de ses 10 premières années
\end{itemize}
\src{Hosseini, Schrimpf, Zhang, Bowman, Zaslavsky \& Fedorenko, Neurobiology of Language 2024}
\end{frame}

\begin{frame}{Piloter le cerveau avec un modèle}
\begin{itemize}
    \item Un modèle de type GPT ajusté sur les réponses fMRI à 1\,000 phrases
    \item Il sert à \textbf{choisir de nouvelles phrases} qui augmentent ou suppriment l'activité du réseau du langage
    \item L'expérience confirme les prédictions chez de nouveaux participants
\end{itemize}
\vspace{0.5em}
\begin{block}{Pourquoi c'est important}
    On passe de la corrélation à un contrôle de l'activité cérébrale.
\end{block}
\src{Tuckute et al., Nature Human Behaviour 2024}
\end{frame}

\begin{frame}{Faire mieux avec le cerveau comme guide}
\begin{itemize}
    \item Utiliser des enregistrements cérébraux comme signal pour améliorer les modèles de langage
    \item Un modèle entraîné à mieux comprendre les récits devient plus aligné au cerveau
\end{itemize}
% TODO : vérifier et choisir 1 ou 2 articles de Toneva (Toneva \& Wehbe 2019 ; Aw \& Toneva 2023)
\src{Toneva \& Wehbe, NeurIPS 2019 ; Aw \& Toneva, ICLR 2023}
\end{frame}

\begin{frame}{L'instruction-tuning aligne-t-il au cerveau ?}
\begin{itemize}
    \item 25 modèles comparés (bruts et ajustés par instructions), 3 jeux de données
    \item L'instruction-tuning augmente l'alignement cérébral d'environ \textbf{6\%} en moyenne
    \item L'alignement suit la taille du modèle et ses connaissances du monde
    \item Mais il n'améliore pas l'alignement avec les temps de lecture
\end{itemize}
\src{Aw, Montariol, AlKhamissi, Schrimpf \& Bosselut, COLM 2024}
\end{frame}

\begin{frame}{L'alignement ne croît pas indéfiniment}
\begin{itemize}
    \item Huit modèles Pythia suivis tout au long de leur entraînement, sur cinq jeux de données cérébrales
    \item L'alignement monte vite, atteint un pic vers 8 milliards de mots, puis stagne
    \item Il suit la compétence \textbf{formelle} (règles du langage) plus que la compétence fonctionnelle
\end{itemize}
\src{AlKhamissi, Tuckute, Tang, Binhuraib, Bosselut \& Schrimpf, EMNLP 2025}
\end{frame}

\begin{frame}{Langage et pensée ne sont pas la même chose}
\begin{itemize}
    \item Compétence \textbf{formelle} : grammaire, bonnes phrases (les LLM excellent)
    \item Compétence \textbf{fonctionnelle} : raisonner, planifier, comprendre le monde (plus fragile)
    \item Chez l'humain, ces deux compétences reposent sur des réseaux cérébraux différents
\end{itemize}
% TODO : vérifier la référence (Mahowald et al., Trends in Cognitive Sciences 2024)
\src{Mahowald, Ivanova, Blank, Kanwisher, Tenenbaum \& Fedorenko, 2024}
\end{frame}

    \section{Q2 : expliquent-ils le comportement humain ?}
\frame{\sectionpage}

\begin{frame}{La surprise d'un mot}
\[
    s(w_t) = -\log p(w_t \mid w_1,\dots,w_{t-1})
\]
\begin{itemize}
    \item Un mot improbable est plus « surprenant »
    \item Les humains lisent plus lentement les mots surprenants
    \item Un modèle de langage fournit directement cette probabilité
\end{itemize}
\src{Levy, Cognition 2008 ; Smith \& Levy, Cognition 2013}
\end{frame}

\begin{frame}{Pourquoi les humains oublient parfois}
\begin{itemize}
    \item La mémoire est limitée : on ne garde pas tout le contexte
    \item Théorie de la \textbf{surprise à contexte dégradé} (lossy-context surprisal)
    \item Explique pourquoi les langues rapprochent les mots liés syntaxiquement
\end{itemize}
% TODO : vérifier les références de Futrell (2015 PNAS dependency length ; 2020 Cognitive Science)
\src{Futrell, Gibson \& Levy, Cognitive Science 2020}
\end{frame}

\begin{frame}{La surprise explique le N400}
\begin{itemize}
    \item Le \textbf{N400} : une onde EEG ($\sim$400 ms) qui augmente quand un mot est inattendu
    \item La surprise de modèles de langage explique de nombreux effets du N400
    \item Reste des cas où l'humain et le modèle divergent
\end{itemize}
% IMAGE : exemple d'onde N400 (à ajouter)
% TODO : vérifier la référence (Michaelov et al., Neurobiology of Language 2024)
\src{Michaelov, Bardolph, Van Petten, Bergen \& Coulson, 2024}
\end{frame}

\begin{frame}{Un meilleur modèle est-il plus humain ?}
\begin{itemize}
    \item Hypothèse admise en anglais : plus la perplexité est basse, plus le modèle prédit les temps de lecture
    \item En japonais, cette règle ne tient pas : \textbf{une perplexité plus basse n'est pas toujours plus humaine}
    \item Il faut évaluer les modèles sur plusieurs langues
\end{itemize}
\begin{block}{Message}
    Plus gros n'est pas automatiquement plus proche de l'humain.
\end{block}
\src{Kuribayashi, Oseki, Ito, Yoshida, Asahara \& Inui, ACL 2021}
\end{frame}

\begin{frame}{Comment le sens se compose}
\begin{itemize}
    \item « maison de chien » $\neq$ « chien de maison »
    \item Des modèles simples de composition sémantique sont testés sur des données fMRI
    \item Une composition apparaît dans le gyrus frontal inférieur gauche, avec un modèle explicite de la combinaison
\end{itemize}
\src{Ciapparelli, Marelli, Graves \& Reverberi, Cerebral Cortex 2025}
\end{frame}

    \section{Q3 : peut-on ouvrir la boîte noire ?}

\begin{frame}{Tester la syntaxe : l'accord sujet-verbe}
\begin{block}{Exemple}
    « Les \textbf{clés} de l'armoire \textbf{sont} sur la table. » \\
    Le verbe s'accorde avec « clés », pas avec « armoire ».
\end{block}
\begin{itemize}
    \item Les LSTM apprennent cet accord à longue distance
    \item Mais ont-ils compris la structure, ou une astuce statistique ?
\end{itemize}
\src{Linzen, Dupoux \& Goldberg, TACL 2016}
\end{frame}

\begin{frame}{Des neurones « nombre » dans un LSTM}
\begin{itemize}
    \item Seulement \textbf{deux neurones} portent l'information singulier/pluriel à longue distance
    \item Un petit circuit « syntaxe » les met à jour au bon moment
    \item On peut intervenir sur ces neurones et voir le modèle se tromper
\end{itemize}
% IMAGE : figure des unités de nombre (à ajouter, avec autorisation)
\src{Lakretz, Kruszewski, Desbordes, Hupkes, Dehaene \& Baroni, NAACL 2019}
\end{frame}

\begin{frame}{Phrases imbriquées : l'humain et le modèle}
\begin{block}{Exemple}
    « Le garçon \textbf{que} les filles \textbf{que} je vois \textbf{aiment} \textbf{court}. »
\end{block}
\begin{itemize}
    \item Difficile pour l'humain comme pour le réseau
    \item Les types d'erreurs se ressemblent sur certains points, pas sur tous
    \item Les transformers ne traitent pas ces phrases comme les humains
\end{itemize}
% TODO : vérifier les deux références (Cognition 2021 ; COLING 2022)
\src{Lakretz et al., Cognition 2021 ; Lakretz, Desbordes, Hupkes \& Dehaene, COLING 2022}
\end{frame}

\begin{frame}{Dans le cerveau : des neurones uniques}
\begin{itemize}
    \item 21 patients implantés, plus de 1\,000 neurones, 1\,400 microfils, 1\,500 contacts iEEG
    \item Phrases lues et entendues
    \item Neurones sensibles aux \textbf{phonèmes} (temporal supérieur) et aux \textbf{lettres} (occipito-temporal ventral)
    \item Neurones \textbf{amodaux}, indépendants de la modalité, dans le temporal moyen et le frontal inférieur
\end{itemize}
\src{Lakretz, Friedmann, King, Mankin, Rangel, Tankus, Dehaene \& Fried, preprint 2024}
\end{frame}

\begin{frame}{Les mêmes outils des deux côtés}
\begin{itemize}
    \item Chez l'humain : un \textbf{localisateur} identifie les régions du langage
    \item Dans un LLM : on identifie de la même façon les unités sélectives au langage, puis on les \textbf{supprime}
    \item Effet causal : leur suppression dégrade les performances langagières
\end{itemize}
% TODO : vérifier la référence (AlKhamissi, Tuckute, Bosselut, Schrimpf, 2025)
\src{AlKhamissi, Tuckute, Bosselut \& Schrimpf, NAACL 2025}
\end{frame}

    \section{Contribution : contraintes cognitives et alignement}

\begin{frame}{Laboratoire}
\begin{itemize}
    \item \textbf{LSCP} : Laboratoire de Sciences Cognitives et Psycholinguistique
    \item ENS-PSL \& CNRS, Paris
\end{itemize}
\vfill
\centering
\includegraphics[height=1cm]{images/lscp.png}\hspace{1cm}\includegraphics[height=1.3cm]{images/ens_psl.png}\hspace{1cm}\includegraphics[height=1cm]{images/cnrs.png}
\end{frame}

\begin{frame}{Question de recherche}
\begin{block}{Les contraintes cognitives rapprochent-elles les modèles de l'humain ?}
    Deux contraintes évaluées : une \textbf{mémoire limitée} et un entraînement sur \textbf{peu de données} (10 à 100 millions de mots).
\end{block}
\begin{itemize}
    \item 118 modèles de langage comparés à des temps de lecture et à des enregistrements cérébraux
    \item Cerveau : fMRI, EEG, MEG et intracrânien
    \item Méthode : analyse de similarité représentationnelle (RSA)
\end{itemize}
\end{frame}

\begin{frame}{Résultats}
\begin{itemize}
    \item Les contraintes augmentent la similarité avec les \textbf{temps de lecture} et la diminuent avec les \textbf{réponses cérébrales}
    \item Les modèles de grande taille présentent le profil inverse
    \item Aucun modèle n'est le meilleur sur les deux cibles
\end{itemize}
\vspace{0.4em}
\begin{block}{Conclusion}
    La similarité comportementale et la similarité cérébrale sont dissociables.
\end{block}
% IMAGE : figure 2 de l'article (comportement vs cerveau), à ajouter
\src{Travail en cours}
\end{frame}

    \section{Bilan et perspectives}

\begin{frame}{État des résultats}
\centering\small
\begin{tabular}{lp{8cm}}
    \textcolor{navy}{\textbf{Domaine}} & \textcolor{navy}{\textbf{Statut}} \\
    \hline
    Alignement modèle-cerveau & Robuste (fMRI, ECoG, MEG), mécanisme débattu \\
    Comportement & Partiel : temps de lecture et N400, avec des limites \\
    Mécanismes & Exemples précis, peu généralisés \\
    Supervision cérébrale & Preuves de principe \\
\end{tabular}

\vspace{0.4em}
{\footnotesize Statuts : estimation de l'auteur d'après les exemples présentés.}
\end{frame}

\begin{frame}{Limites}
\begin{itemize}
    \item \textbf{Corrélation n'est pas mécanisme} : un score d'alignement élevé n'établit pas l'identité du calcul
    \item Des scores élevés peuvent s'expliquer par des facteurs simples (longueur, position dans la phrase)
    \item Données majoritairement en anglais
    \item Données intracrâniennes rares, recueillies chez des patients
\end{itemize}
% TODO : ajouter la référence exacte pour le deuxième point
\end{frame}

\begin{frame}{Conclusions}
\begin{enumerate}
    \item Les modèles de langage \textbf{prédisent} l'activité du réseau du langage
    \item Ils \textbf{expliquent} une partie du comportement humain, mais pas tout
    \item Leurs représentations internes peuvent être \textbf{disséquées} avec les mêmes questions que celles posées au cerveau
\end{enumerate}
\end{frame}

\section*{}
\begin{frame}{}
    \centering
    \Huge\bfseries
    \textcolor{navy}{Merci}\\[0.5em]
    \Large Questions ?
\end{frame}


\end{document}
